% Generated by tex/build_swebok_structure.py from tex/chapters/ch01/06_要求仕様化.tex
\subsubsection{構造化自然言語による要求仕様化}

\term{構造化自然言語}とは、自然言語を使いながら、書き方に制約を加える方法である。代表例は、\term{主体・動作形式}である。「いつ」「誰が」「何をする」「どの条件で」という形にそろえると、曖昧さを減らせる。

\MiniBox{例}{「注文が出荷されたとき、システムは請求書を作成する。ただし、注文条件が前払いの場合を除く。」この文では、きっかけ、主体、動作、例外条件が明示されている。}

他にも、\term{ユースケース仕様}、\term{ユーザーストーリー}、\term{決定表}などが構造化自然言語の例である。ユーザーストーリーは、「ある役割として、ある能力がほしい。なぜなら、ある便益を得たいからである」という形式で、利用者視点の価値を表す。
